% GENERATED FILE. Edit canonical hl-activity contracts, then run npm run generate:books.
% canonical-activities: 21

\chapter*{Review Questions}
\addcontentsline{toc}{chapter}{Review Questions}
\markboth{Review Questions}{Review Questions}

Try these without looking ahead. Every prompt comes from the same canonical lesson data as its chapter. Answers begin in the next section.

\section*{Chapter 6}

\noindent\begin{minipage}[t]{\linewidth}
\raggedright
\hypertarget{review-PA-C06-panj-convergence-borrowing}{\textbf{6.1}} \emph{Why two branches both arrived at \emph{panj}}\par
\small Was Punjabi panj borrowed from Persian?
\end{minipage}\par\medskip

\section*{Chapter 7}

\noindent\begin{minipage}[t]{\linewidth}
\raggedright
\hypertarget{review-PA-C07-janna-two-js}{\textbf{7.1}} \emph{\pa{ਜਾਣਨਾ} — “to know,” one letter from “to go”}\par
\small Which Punjabi verb means to know, jana or janna?
\end{minipage}\par\medskip

\section*{Chapter 8}

\noindent\begin{minipage}[t]{\linewidth}
\raggedright
\hypertarget{review-PA-C08-sochna-soch}{\textbf{8.1}} \emph{\pa{ਸੋਚਣਾ} — “to think,” and it once meant “to burn”}\par
\small The Punjabi noun soch means a thought and one other thing. What?
\end{minipage}\par\medskip

\noindent\begin{minipage}[t]{\linewidth}
\raggedright
\hypertarget{review-PA-C08-samajhna-tone}{\textbf{8.2}} \emph{\pa{ਸਮਝਣਾ} — “to understand,” and the tone you can hear in it}\par
\small In samajhna the jh is not heard as breath. What is heard instead?
\end{minipage}\par\medskip

\noindent\begin{minipage}[t]{\linewidth}
\raggedright
\hypertarget{review-PA-C08-parhna-subjoined}{\textbf{8.3}} \emph{\pa{ਪੜ੍ਹਨਾ} — “to read,” which used to mean “to say out loud”}\par
\small In the cluster rrha, where is the h written relative to the letter before it?
\end{minipage}\par\medskip

\noindent\begin{minipage}[t]{\linewidth}
\raggedright
\hypertarget{review-PA-C08-likhna-four-roots}{\textbf{8.4}} \emph{\pa{ਲਿਖਣਾ} — “to write,” which is scratching}\par
\small What did the Sanskrit root behind likhna mean before it meant to write?
\end{minipage}\par\medskip

\section*{Chapter 9}

\noindent\begin{minipage}[t]{\linewidth}
\raggedright
\hypertarget{review-PA-C09-laina-labhate}{\textbf{9.1}} \emph{\pa{ਲੈਣਾ} — “to take,” worn down to almost nothing}\par
\small Which Sanskrit verb is laina worn down from?
\end{minipage}\par\medskip

\noindent\begin{minipage}[t]{\linewidth}
\raggedright
\hypertarget{review-PA-C09-puchhna-addak}{\textbf{9.2}} \emph{\pa{ਪੁੱਛਣਾ} — “to ask,” and a Persian cousin that is not a loan}\par
\small What does the addak mark tell you to do to the consonant after it?
\end{minipage}\par\medskip

\noindent\begin{minipage}[t]{\linewidth}
\raggedright
\hypertarget{review-PA-C09-madad-karna-root}{\textbf{9.3}} \emph{\pa{ਮਦਦ} \pa{ਕਰਨਾ} — “to help,” which is two words}\par
\small Which language does the word madad come from originally?
\end{minipage}\par\medskip

\noindent\begin{minipage}[t]{\linewidth}
\raggedright
\hypertarget{review-PA-C09-pasand-dative}{\textbf{9.4}} \emph{\pa{ਪਸੰਦ} — Punjabi does not like things; things please Punjabi}\par
\small Type the Punjabi for 'to me' — the form main takes when it cannot be the subject.
\end{minipage}\par\medskip

\section*{Chapter 10}

\noindent\begin{minipage}[t]{\linewidth}
\raggedright
\hypertarget{review-PA-C10-paani-please}{\textbf{10.1}} \emph{\pa{ਪਾਣੀ} — the first thing you'll actually ask for}\par
\small Ask for water politely, in Punjabi.
\end{minipage}\par\medskip

\noindent\begin{minipage}[t]{\linewidth}
\raggedright
\hypertarget{review-PA-C10-chaa-tone}{\textbf{10.2}} \emph{\pa{ਚਾਹ} — the \pa{ਹ} that is not a consonant}\par
\small Is chaa's tone falling or level?
\end{minipage}\par\medskip

\noindent\begin{minipage}[t]{\linewidth}
\raggedright
\hypertarget{review-PA-C10-roti-requests}{\textbf{10.3}} \emph{\pa{ਰੋਟੀ} — a request pattern, run four times}\par
\small Ask for bread politely, in Punjabi.
\end{minipage}\par\medskip

\section*{Chapter 11}

\noindent\begin{minipage}[t]{\linewidth}
\raggedright
\hypertarget{review-PA-C11-parivar-split}{\textbf{11.1}} \emph{\pa{ਪਰਿਵਾਰ} — what surrounds you}\par
\small parivar splits into two Sanskrit pieces meaning 'around' and 'to cover'. What is the whole word's literal sense?
\end{minipage}\par\medskip

\noindent\begin{minipage}[t]{\linewidth}
\raggedright
\hypertarget{review-PA-C11-bhara-tone}{\textbf{11.2}} \emph{\pa{ਭਰਾ} — a breathy letter you have not met yet}\par
\small bharā's bh is said as which consonant, and what tone does the vowel after it carry?
\end{minipage}\par\medskip

\noindent\begin{minipage}[t]{\linewidth}
\raggedright
\hypertarget{review-PA-C11-bhain-cognate}{\textbf{11.3}} \emph{\pa{ਭੈਣ} — the same tone, a different family tree}\par
\small Is bhaiṇ the same word as English sister by unbroken descent, the way bharā is the same word as brother?
\end{minipage}\par\medskip

\section*{Chapter 12}

\noindent\begin{minipage}[t]{\linewidth}
\raggedright
\hypertarget{review-PA-C12-kann-false-friend}{\textbf{12.1}} \emph{\pa{ਕੰਨ} — a false friend hiding in plain sight}\par
\small Does kann, 'ear', share a root with karnā, 'to do'?
\end{minipage}\par\medskip

\noindent\begin{minipage}[t]{\linewidth}
\raggedright
\hypertarget{review-PA-C12-munh-tone}{\textbf{12.2}} \emph{\pa{ਮੂੰਹ} — a consonant that wore all the way down to a tone}\par
\small What Sanskrit word does mūnh come from, and what learned Punjabi word keeps the same source intact?
\end{minipage}\par\medskip

\noindent\begin{minipage}[t]{\linewidth}
\raggedright
\hypertarget{review-PA-C12-nakk-face}{\textbf{12.3}} \emph{\pa{ਨੱਕ} — the whole face, run back}\par
\small Name all four Punjabi face words taught in this chapter, in any order.
\end{minipage}\par\medskip

\section*{Chapter 13}

\noindent\begin{minipage}[t]{\linewidth}
\raggedright
\hypertarget{review-PA-C13-dil-source}{\textbf{13.1}} \emph{\pa{ਦਿਲ} — the one body word that is not Sanskrit}\par
\small Is dil, Punjabi's everyday word for heart, inherited from Sanskrit or borrowed from Persian?
\end{minipage}\par\medskip

\noindent\begin{minipage}[t]{\linewidth}
\raggedright
\hypertarget{review-PA-C13-sir-cousin}{\textbf{13.2}} \emph{\pa{ਸਿਰ} — a cousin of \textquotedblleft{}horn,\textquotedblright{} and a twin nobody borrowed}\par
\small sir shares its Proto-Indo-European root with which English word — head, or horn?
\end{minipage}\par\medskip

\chapter*{Answer Key}
\addcontentsline{toc}{chapter}{Answer Key}
\markboth{Answer Key}{Answer Key}

The first response is the canonical display answer. When a question accepts other authored forms, they appear underneath.

\section*{Chapter 6}

\noindent\begin{minipage}[t]{\linewidth}
\raggedright
\hyperlink{review-PA-C06-panj-convergence-borrowing}{\textbf{6.1}} \emph{Why two branches both arrived at \emph{panj}}\par
\small \textbf{Answer:} no\par
\footnotesize \textbf{Also accepted:} no it was not; it inherits Sanskrit pañca; it comes from Sanskrit pañca
\end{minipage}\par\medskip

\section*{Chapter 7}

\noindent\begin{minipage}[t]{\linewidth}
\raggedright
\hyperlink{review-PA-C07-janna-two-js}{\textbf{7.1}} \emph{\pa{ਜਾਣਨਾ} — “to know,” one letter from “to go”}\par
\small \textbf{Answer:} janna\par
\footnotesize \textbf{Also accepted:} jāṇnā; jannaa; jaanna
\end{minipage}\par\medskip

\section*{Chapter 8}

\noindent\begin{minipage}[t]{\linewidth}
\raggedright
\hyperlink{review-PA-C08-sochna-soch}{\textbf{8.1}} \emph{\pa{ਸੋਚਣਾ} — “to think,” and it once meant “to burn”}\par
\small \textbf{Answer:} a worry\par
\footnotesize \textbf{Also accepted:} worry; worrying; anxiety; a worry or a care
\end{minipage}\par\medskip

\noindent\begin{minipage}[t]{\linewidth}
\raggedright
\hyperlink{review-PA-C08-samajhna-tone}{\textbf{8.2}} \emph{\pa{ਸਮਝਣਾ} — “to understand,” and the tone you can hear in it}\par
\small \textbf{Answer:} a falling tone\par
\footnotesize \textbf{Also accepted:} falling tone; a falling pitch; the pitch falls; high falling tone; tone
\end{minipage}\par\medskip

\noindent\begin{minipage}[t]{\linewidth}
\raggedright
\hyperlink{review-PA-C08-parhna-subjoined}{\textbf{8.3}} \emph{\pa{ਪੜ੍ਹਨਾ} — “to read,” which used to mean “to say out loud”}\par
\small \textbf{Answer:} underneath\par
\footnotesize \textbf{Also accepted:} below; under; beneath; underneath it; under it
\end{minipage}\par\medskip

\noindent\begin{minipage}[t]{\linewidth}
\raggedright
\hyperlink{review-PA-C08-likhna-four-roots}{\textbf{8.4}} \emph{\pa{ਲਿਖਣਾ} — “to write,” which is scratching}\par
\small \textbf{Answer:} to scratch\par
\footnotesize \textbf{Also accepted:} scratch; scrape; to scrape; scratching; scraping
\end{minipage}\par\medskip

\section*{Chapter 9}

\noindent\begin{minipage}[t]{\linewidth}
\raggedright
\hyperlink{review-PA-C09-laina-labhate}{\textbf{9.1}} \emph{\pa{ਲੈਣਾ} — “to take,” worn down to almost nothing}\par
\small \textbf{Answer:} labhate\par
\footnotesize \textbf{Also accepted:} labh; labh-; labhati; sanskrit labhate
\end{minipage}\par\medskip

\noindent\begin{minipage}[t]{\linewidth}
\raggedright
\hyperlink{review-PA-C09-puchhna-addak}{\textbf{9.2}} \emph{\pa{ਪੁੱਛਣਾ} — “to ask,” and a Persian cousin that is not a loan}\par
\small \textbf{Answer:} double it\par
\footnotesize \textbf{Also accepted:} double; doubles it; hold it long; lengthen it
\end{minipage}\par\medskip

\noindent\begin{minipage}[t]{\linewidth}
\raggedright
\hyperlink{review-PA-C09-madad-karna-root}{\textbf{9.3}} \emph{\pa{ਮਦਦ} \pa{ਕਰਨਾ} — “to help,” which is two words}\par
\small \textbf{Answer:} Arabic\par
\footnotesize \textbf{Also accepted:} from arabic; arabic via persian; arabic through persian
\end{minipage}\par\medskip

\noindent\begin{minipage}[t]{\linewidth}
\raggedright
\hyperlink{review-PA-C09-pasand-dative}{\textbf{9.4}} \emph{\pa{ਪਸੰਦ} — Punjabi does not like things; things please Punjabi}\par
\small \textbf{Answer:} \pa{ਮੈਨੂੰ}\par
\footnotesize \textbf{Also accepted:} mainu; mainun; mainoo; mainū̃
\end{minipage}\par\medskip

\section*{Chapter 10}

\noindent\begin{minipage}[t]{\linewidth}
\raggedright
\hyperlink{review-PA-C10-paani-please}{\textbf{10.1}} \emph{\pa{ਪਾਣੀ} — the first thing you'll actually ask for}\par
\small \textbf{Answer:} \pa{ਪਾਣੀ}, \pa{ਕਿਰਪਾ} \pa{ਕਰਕੇ}\par
\footnotesize \textbf{Also accepted:} pani kirpa karke; paani, kirpa karke; pāṇī, kirpā karke; pani kirpa karke.
\end{minipage}\par\medskip

\noindent\begin{minipage}[t]{\linewidth}
\raggedright
\hyperlink{review-PA-C10-chaa-tone}{\textbf{10.2}} \emph{\pa{ਚਾਹ} — the \pa{ਹ} that is not a consonant}\par
\small \textbf{Answer:} level\par
\footnotesize \textbf{Also accepted:} high level; level high; high and level; stays high
\end{minipage}\par\medskip

\noindent\begin{minipage}[t]{\linewidth}
\raggedright
\hyperlink{review-PA-C10-roti-requests}{\textbf{10.3}} \emph{\pa{ਰੋਟੀ} — a request pattern, run four times}\par
\small \textbf{Answer:} \pa{ਰੋਟੀ}, \pa{ਕਿਰਪਾ} \pa{ਕਰਕੇ}\par
\footnotesize \textbf{Also accepted:} roti kirpa karke; roṭī, kirpā karke; roti, kirpa karke.
\end{minipage}\par\medskip

\section*{Chapter 11}

\noindent\begin{minipage}[t]{\linewidth}
\raggedright
\hyperlink{review-PA-C11-parivar-split}{\textbf{11.1}} \emph{\pa{ਪਰਿਵਾਰ} — what surrounds you}\par
\small \textbf{Answer:} what surrounds you\par
\footnotesize \textbf{Also accepted:} that which surrounds; surrounding; what encloses you; that which covers around
\end{minipage}\par\medskip

\noindent\begin{minipage}[t]{\linewidth}
\raggedright
\hyperlink{review-PA-C11-bhara-tone}{\textbf{11.2}} \emph{\pa{ਭਰਾ} — a breathy letter you have not met yet}\par
\small \textbf{Answer:} p, low tone\par
\footnotesize \textbf{Also accepted:} p and low; plain p, low; p, low; voiceless p and a low tone
\end{minipage}\par\medskip

\noindent\begin{minipage}[t]{\linewidth}
\raggedright
\hyperlink{review-PA-C11-bhain-cognate}{\textbf{11.3}} \emph{\pa{ਭੈਣ} — the same tone, a different family tree}\par
\small \textbf{Answer:} no\par
\footnotesize \textbf{Also accepted:} no it is not; no, unrelated; no different root
\end{minipage}\par\medskip

\section*{Chapter 12}

\noindent\begin{minipage}[t]{\linewidth}
\raggedright
\hyperlink{review-PA-C12-kann-false-friend}{\textbf{12.1}} \emph{\pa{ਕੰਨ} — a false friend hiding in plain sight}\par
\small \textbf{Answer:} no\par
\footnotesize \textbf{Also accepted:} no it does not; no unrelated; no, coincidence
\end{minipage}\par\medskip

\noindent\begin{minipage}[t]{\linewidth}
\raggedright
\hyperlink{review-PA-C12-munh-tone}{\textbf{12.2}} \emph{\pa{ਮੂੰਹ} — a consonant that wore all the way down to a tone}\par
\small \textbf{Answer:} mukha, and mukh\par
\footnotesize \textbf{Also accepted:} mukha and mukh; mukha; mukh; mukha, mukh
\end{minipage}\par\medskip

\noindent\begin{minipage}[t]{\linewidth}
\raggedright
\hyperlink{review-PA-C12-nakk-face}{\textbf{12.3}} \emph{\pa{ਨੱਕ} — the whole face, run back}\par
\small \textbf{Answer:} akkh, kann, munh, nakk\par
\footnotesize \textbf{Also accepted:} akkh kann munh nakk; \pa{ਅੱਖ} \pa{ਕੰਨ} \pa{ਮੂੰਹ} \pa{ਨੱਕ}; eye ear mouth nose
\end{minipage}\par\medskip

\section*{Chapter 13}

\noindent\begin{minipage}[t]{\linewidth}
\raggedright
\hyperlink{review-PA-C13-dil-source}{\textbf{13.1}} \emph{\pa{ਦਿਲ} — the one body word that is not Sanskrit}\par
\small \textbf{Answer:} borrowed from Persian\par
\footnotesize \textbf{Also accepted:} persian; borrowed persian; persian loan
\end{minipage}\par\medskip

\noindent\begin{minipage}[t]{\linewidth}
\raggedright
\hyperlink{review-PA-C13-sir-cousin}{\textbf{13.2}} \emph{\pa{ਸਿਰ} — a cousin of \textquotedblleft{}horn,\textquotedblright{} and a twin nobody borrowed}\par
\small \textbf{Answer:} horn\par
\footnotesize \textbf{Also accepted:} horn not head; its horn
\end{minipage}\par\medskip
